\scanpage{48}
puisse espérer dans cette direction est ceci :

a) Si $n$ \inserted{$\geq 3$} n'est pas premier $\ne 2$ ni de la forme $2p$,
$p$ premier, trouver un $\mathrm{SR}n\mathrm S$ sur $A=\mathbb Z[\frac1n]$

b) Si $n$ est premier impair ou $n=4$,
trouver un $\mathrm{SR}n\mathrm S$ sur $A=\mathbb Z$

c) Si $n=2p$, $p$ premier impair, trouver
un $\mathrm{SR}n\mathrm S$ sur $A=\mathbb Z[\frac12]$.

La solution à a) est évidente, en prenant
\[
\left\{\begin{aligned}
Z_0&=\unsure{\mu_{n\,A}|\mu_{n_A}|\mu_{n,A}}\subset\mathrm{Sl}(2)_{\mathbb Z}\subset\underline{\mathrm{Aff}}\,\mathbb E^2_{/A}\\
\sigma_0&=\begin{pmatrix}0&1\\1&0\end{pmatrix}\in\mathrm{Sl}(2,\mathbb Z)\subset\mathrm{Aff}(\mathbb E_A^2)\\
s_0&=(1,1)\in\mathbb E^2(A)^\sigma\qquad S_0=Z.s_0
\end{aligned}\right.
\]
On trouve le \deleted{\unsure{grou|genre|groupe}} $\mathrm{SR}n\mathrm S$ type $(S_0,E_0)$
dont le schéma des autom est $G=$
$\{\pm1\}_A.\mu_{n\,A}$ (produit $\frac12$ direct) -- la
cat. fibrée des $\mathrm{SR}n\mathrm S$ s'identifie alors
à celle des $G$-torseurs. Ceci
marche encore pour \uline{tt} $n\geq3$, pourvu
qu'on se place sur $\mathbb Z[1/n]$ --
donc dans les cas b) et c), on \uline{perd}
la car. $p$, où les choses sont spéciales.
Bien sûr, \deleted{sur les \unsure{co|corps}} \inserted{dans ces cas-là}, sur le corps
% archive sheet number: 47
